% Self-contained research appendix fragment; amsmath and amssymb suffice.
% The main obstruction is model-theoretic, not a source-size lower bound.
\section{A finite-parameter obstruction to removing numerical cases}
\label{app:finite-parameter-obstruction}

\paragraph{Antecedents and precise contribution.}
The finite-parameter involution versus definable-linear-order method,
including the treatment of quotient representatives, is already explicit
in the nested unary-arithmetic amalgam drafts~\cite{prioramalgam}.
The simpler product-model independence of coverage is also already
present~\cite{priorproduct}. The construction below extends this line
of argument using a countable eventually constant power for full-arity
arithmetic, all true universal Horn laws, and arbitrary finitary
expansions with atomic or conjunctive induction. No priority for the
automorphism method or the elementary product observation is claimed.

This section proves a stronger limitation than failure of one proposed
term compiler.  A countable model built from finitary objects obstructs
\emph{any} ordinary definable bi-interpretation with iPRA, even when
finite tuples of parameters, different interpretation domains,
finite-dimensional coding, and definable quotients are allowed.
It applies to the ordinary first-order realization of the numerical
theory left after removing constructor coverage.  The scope concerning
native calculi with unusual logical connectives is stated explicitly
below.

\subsection{The ordinary theory under consideration}

Use the usual one-sorted first-order signature with a function symbol
for every fixed primitive-recursive description.  Different descriptions
may name the same standard function.  Let $T_{\rm alg}$ have ordinary
intuitionistic first-order logic with equality, all PR defining equations,
constructor separation, and induction for each literal equation
$E(n,\vec p)$:
\[
 \forall\vec p\,
 \left[
 \left(E(0,\vec p)\land
       \forall n\,(E(n,\vec p)\to E(Sn,\vec p))\right)
 \to\forall n\,E(n,\vec p)
 \right].
 \tag{EI}
\]
One may take separation to include both $0\ne Sx$ and successor
injectivity; the construction below satisfies them.  Equivalently,
the weaker boundary axiom $0\ne1$ may be used when the others are
derived from the chosen PR descriptions.  Constructor coverage is
deliberately absent.

Although the theory is intuitionistic, it has ordinary classical
structures as one-world models.  Every interpretation or definable
isomorphism proved intuitionistically must remain valid in these
structures.  They therefore suffice for the obstruction.

\subsection{A countable model of all literal-equation induction}

Let $\mathcal E$ have domain
\[
 \left\{a:\mathbb N\to\mathbb N:
   \text{there exist }c\in\mathbb N\text{ and finite }F\subseteq\mathbb N
   \text{ with }a(i)=c\text{ for all }i\notin F\right\}.
 \tag{Domain}
\]
These are the eventually constant sequences.  Each element has a
finite description consisting of its eventual value and its finitely
many exceptional coordinate values, so the domain is countable.
Interpret every PR symbol pointwise using its ordinary standard
numerical meaning.  Constants are constant sequences; in particular
\[
 0(i)=0,\qquad S(a)(i)=a(i)+1.
\]

\paragraph{Closure and arithmetic equations.}
If $a_1,\ldots,a_r$ are eventually constant, then outside the union
of their finite exception sets every tuple of their coordinate values
is the same.  Applying any fixed $r$-ary PR function consequently gives
another eventually constant sequence.  Every defining equation is true
coordinatewise.  Successor is injective and its range omits zero,
also coordinatewise.  Equality in this structure is actual equality of
sequences.

\paragraph{Lemma: every instance of EI holds in $\mathcal E$.}
Fix an equation $E(n,\vec p)$ and an arbitrary finite tuple of parameter
sequences $\vec p$.  For each coordinate $i$, put
\[
 U_i=\{m\in\mathbb N:
       E(m,\vec p(i))\text{ holds in the standard natural numbers}\}.
\]
Term evaluation is pointwise, so
\[
 \mathcal E\models E(a,\vec p)
 \quad\Longleftrightarrow\quad
 (\forall i)\ a(i)\in U_i.
 \tag{Coordinate factorization}
\]
Suppose the antecedent of EI holds.  Its base gives $0\in U_i$ for
every $i$.  Fix $i$ and $m\in U_i$, and form the finite-support
sequence
\[
 a_{i,m}(j)=
 \begin{cases}m,&j=i,\\0,&j\ne i.\end{cases}
\]
This element belongs to $\mathcal E$ and satisfies $E(a_{i,m},\vec p)$:
at $i$ this is the choice of $m$, and at every other coordinate it is
the known base case.  The global induction-step premise gives
$E(Sa_{i,m},\vec p)$, hence $m+1\in U_i$.  Thus each $U_i$ contains
zero and is successor closed.  External ordinary induction in the
standard natural numbers gives $U_i=\mathbb N$ for every $i$.
Coordinate factorization now gives $E(a,\vec p)$ for every
$a\in\mathcal E$.  This proves the instance and hence the lemma.
\hfill$\square$

This argument applies to every fixed PR term equation, with arbitrary
finite parameter tuples.  It proves the full axiom schema and therefore
also validates the contextual equation-induction rule under arbitrary
retained assumptions.  It is not merely a verification of induction
for parameter-free equations.

\paragraph{What fails.}
Let $e_i$ be the sequence that is one at $i$ and zero elsewhere.
It is neither zero nor a successor, since every successor sequence has
a positive value at every coordinate.  Therefore constructor coverage
fails at $e_i$.

There is also a direct failed quantifier-free induction instance.  Let
$\operatorname{sg}$ be the usual PR sign function, zero at zero and
one at successors, and let
\[
 \varphi(n):=
   (\operatorname{sg}(n)=0)\lor(\operatorname{sg}(n)=1).
\]
Its base holds, and $\varphi(Sa)$ holds for every $a$, since
$\operatorname{sg}(Sa)$ is the constant-one sequence.  But
$\operatorname{sg}(e_i)=e_i$ is neither constant zero nor constant one,
so $\varphi(e_i)$ fails.  The missing feature is a disjunctive range
assertion, not an inability to perform numerical computation.

For precision, the standard PR equality characteristic still reflects
actual equality in this model: its value is the zero sequence exactly
when its two arguments are equal at every coordinate.  What fails is
the assertion that this characteristic is always \emph{globally}
equal to zero or globally equal to one.  These are different claims.

\subsection{Finite parameters cannot remove the symmetry}

\paragraph{Lemma.}
For every finite tuple $\vec a$ of elements of $\mathcal E$, the
parameter expansion $(\mathcal E,\vec a)$ has a nonidentity
automorphism of order two.

\paragraph{Proof.}
Let $F$ be the union of finite exception sets for all entries of
$\vec a$.  Choose distinct coordinates $i,j\notin F$.  Exchange those
two coordinates in every sequence and leave all others fixed.  This
map preserves eventual constancy and commutes with every pointwise PR
operation.  It fixes each entry of $\vec a$, because that entry takes
its eventual value at both $i$ and $j$.  It is nonidentity because it
interchanges $e_i$ and $e_j$, and applying it twice is the identity.
\hfill$\square$

In fact the same construction gives every finite permutation of
coordinates outside $F$, but an order-two automorphism is enough.
It follows already that no strict linear order on $\mathcal E$ can
be defined with a finite tuple of parameters.

\subsection{Why bi-interpretability cannot conceal this symmetry}

Every ordinary model of iPRA has a parameter-free definable strict
linear order, for example
\[
 x<y\quad\Longleftrightarrow\quad (Sx)\mathbin{\dotminus}y=0,
\]
where $\dotminus$ is truncated subtraction.  The usual order laws
are elementary PRA theorems, proved with quantifier-free induction in
the language of PR comparisons.  Thus every automorphism of an iPRA
model preserves a strict linear order.

An order-preserving automorphism $\sigma$ of finite order is the
identity.  If $a<\sigma(a)$, repeated application gives a strict cycle
$a<\sigma(a)<\cdots<\sigma^r(a)=a$.  If $\sigma(a)<a$, the same
argument with reversed inequalities gives a contradiction.  Linearity
therefore gives $\sigma(a)=a$ for every $a$.  Consequently the
automorphism group of every iPRA model has no nonidentity element
of finite order.

Here is the needed consequence of bi-interpretability, with its proof.
Suppose a structure $M$ interprets $N$ by a finite-dimensional
definable domain and a definable equality relation, possibly a quotient.
An automorphism of $M$ fixing the finitely many interpretation
parameters acts coordinatewise on representing tuples and therefore
induces an automorphism of $N$.  This is a group homomorphism.

Suppose, in addition, $N$ interprets a copy of $M$ and the composite
copy is definably isomorphic to $M$, as required for a bi-interpretation.
Collect into one finite tuple $\vec c$ all parameters used by the
first interpretation, the definable composite isomorphism, and finite
representatives in $M$ of parameters used by the second interpretation.
The induced homomorphism
\[
 \operatorname{Aut}(M/\vec c)\longrightarrow\operatorname{Aut}(N)
 \tag{Automorphism action}
\]
is injective.  Indeed, if $\sigma$ induces the identity on $N$, it
induces the identity on the interpreted copy of $M$.  The graph of
the composite isomorphism is definable over parameters fixed by
$\sigma$, so this isomorphism commutes with the two automorphism
actions.  Its injectivity then forces $\sigma$ to be the identity on
$M$.  Only this one composite is needed for the argument.

This proof accommodates finite tuples, definable quotients, different
interpretation domains, and finite parameter data.  It does not assume
that either interpretation has the identity domain or preserves each
element literally.

\paragraph{Theorem.}
The structure $\mathcal E$ is not bi-interpretable with any ordinary
model of iPRA by ordinary definable interpretations with finite
parameter data.  Consequently $T_{\rm alg}$ and iPRA are not
bi-interpretable in that sense, even with finite parameters.

\paragraph{Proof.}
For a hypothetical bi-interpretation, take the finite tuple $\vec c$
in Automorphism action.  The preceding symmetry lemma supplies a
nonidentity element of order two in
$\operatorname{Aut}(\mathcal E/\vec c)$.  Its injective image would
be a nonidentity finite-order automorphism of an iPRA model, which
the definable linear order excludes.  Since $\mathcal E$ is a model
of $T_{\rm alg}$, any theory-level bi-interpretation would apply to it
and give the same contradiction.
\hfill$\square$

\subsection{All true universal Horn laws still do not suffice}

The obstruction persists after adjoining \emph{any collection} of
universal Horn sentences in the fixed PR signature that are true in
the standard natural numbers.  Here a universal Horn sentence has
the form
\[
 \forall\vec x\,
   (E_1(\vec x)\land\cdots\land E_r(\vec x)\to F(\vec x)),
\]
where the premises are equations and the conclusion is an equation
or falsity; the list of premises may be empty.

To verify preservation directly, take any tuple of eventually constant
sequences satisfying all premises.  At each coordinate, the numerical
tuple satisfies all corresponding standard premises.  Truth of the
Horn law in $\mathbb N$ gives its equation conclusion at that
coordinate, and doing this at every coordinate gives actual equality
of the output sequences.  If the conclusion is falsity, even one
coordinate contradicts the assumed standard truth of the Horn law,
so its premises cannot all hold.  No closure-under-products theorem
from model theory is needed for this verification.

Thus $\mathcal E$ satisfies not just the defining equations but every
true universal Horn arithmetic law, while still satisfying EI and
still having the same finite-parameter symmetries.  Even the theory
of all such true Horn laws together with literal-equation induction
cannot be bi-interpretable with iPRA in the stated sense.

Within this fixed numerical signature, adding more true equations
or universal Horn arithmetic axioms therefore cannot repair the
omission of cases.  An additional ingredient must exclude the model;
constructor coverage and Boolean-range disjunctions are concrete
examples that do so.  This is not a prohibition on arbitrary alternative
domains, new operations with different intended meanings, or other
non-Horn axiomatizations.

\paragraph{Generalization to arbitrary finitary expansions.}
The same argument applies to any ordinary finitary first-order expansion
$\mathbb N^*$ of the standard structure $(\mathbb N,0,S)$, not only
to the PR signature.  Interpret each additional function pointwise on
eventually constant sequences, and each relation by requiring its
standard interpretation at every coordinate.  Closure under the
functions follows from eventual constancy; atomic formulas factor
coordinatewise by definition.  The finite-support-sequence proof
therefore validates induction for any atomic formula, and even for
any finite conjunction of atomic formulas.  All universal Horn
sentences true in $\mathbb N^*$ remain true by the same coordinate
argument.  Coordinate transpositions preserve all these functions and
relations and still fix any prescribed finite tuple of parameters.
Consequently such Horn laws plus atomic, or finite-conjunction,
induction cannot force a theory bi-interpretable with iPRA by ordinary
definable interpretations with finite parameter data.

This generalization still starts from an intended standard zero and
successor on $\mathbb N$ and uses ordinary model-theoretic
interpretations.  It is not a theorem about every exotic semantics
or every possible induction mechanism on another kind of object.

\subsection{Exact application to the progressive native calculus}

Interpret the native progressive symbols by
\[
 \mathsf P_yB=\forall y(B(y)\to B(Sy)),\qquad
 \mathsf H_{x,y}B=\exists x\,\forall y(B(x,y)\to B(x,Sy)).
\]
Every progressive-calculus rule other than BoolCases (equivalently,
the derived constructor Cases interface) is sound in
$T_{\rm alg}$: transport, HI, HE, and Step have their ordinary logical
meanings, and equation-only Iter is EI.  Consequently $\mathcal E$
is also a sound standard-semantics model of that case-free rule system.
The theorem rules out rescuing its ordinary first-order realization
by a finite-parameter definable bi-interpretation on a different
domain or quotient.

There is a necessary distinction.  Soundness of a native proof calculus
in an ordinary theory is not by itself a proof that the native calculus
derives all the ordinary logical rules of that theory.  Nor does an
arbitrary proof translation that changes the meanings of logical
connectives automatically give an ordinary model interpretation.
The automorphism theorem applies to standard model-theoretic
bi-interpretability, or to a native proposal whose claimed interpretations
are proved to induce those definable model interpretations under the
displayed semantics.  It is not presented as a theorem about every
conceivable nonstandard meaning of ``proof equivalence.''

The result is also not a lower bound on Python source size or on every
possible axiomatization of iPRA.  It establishes a precise obstruction
to a substantial class of proposed arithmetic simplifications, including
all purely equational or true universal Horn repairs in this signature.
It does not exclude equiconsistency, nor does the automorphism argument
by itself exclude interpretations lacking a definable recovery of the
source model.  The bi-interpretation composite is the essential extra
requirement used by the proof.
